\sectioncentered*{Заключение}
\addcontentsline{toc}{section}{Заключение}

В лабораторной работе реализована линейная рециркуляционная сеть для сжатия и восстановления растровых изображений BMP. Изображение разбивается на блоки, нормализуется и обучается автоассоциативно через скрытый слой размерности $p$. Для варианта с адаптивным коэффициентом обучения и нормированными весами получено сжатие $Z > 1$, показаны исходное и восстановленное изображения, а также график убывания суммарной ошибки по эпохам.

\textbf{Семантическая составляющая.} Сжатие интерпретируется как отображение вектора блока в пространство меньшей размерности с последующим обратным отображением; качество определяется согласованностью матриц $W$ и $W'$ и выбором $p$.

\textbf{Прагматическая составляющая.} Получен воспроизводимый программный стенд на Python с интерактивным запуском, сохранением сжатых представлений, весов, метрик и графиков, пригодный для сравнения режимов обучения и параметров блока.

\begin{thebibliography}{9}
	\addcontentsline{toc}{section}{Список использованных источников}

	\bibitem{method}
	Методические указания к лабораторной работе №3 по дисциплине <<Модели решения задач в интеллектуальных системах>>. \mdash  Минск : БГУИР.

	\bibitem{haykin}
	Хайкин С. Нейронные сети: полный курс. \mdash  2-е изд. \mdash  М. : Вильямс, 2006. \mdash  1104~с.

	\bibitem{numpy}
	NumPy documentation [Электронный ресурс]. \mdash  Режим доступа: \url{https://numpy.org/doc/}. \mdash  Дата доступа: 02.10.2026.

	\bibitem{pillow}
	Pillow (PIL Fork) documentation [Электронный ресурс]. \mdash  Режим доступа: \url{https://pillow.readthedocs.io/}. \mdash  Дата доступа: 02.10.2026.
\end{thebibliography}
